\chapter{2013年山东卷（理）}

\section{单选题}

\begin{problem}
复数 \(z\) 满足 \((z - 3)(2 - \i) = 5\)（\(\i\) 为虚数单位），则 \(z\) 的共轭复数 \(\overline{z}\) 为（\quad）

\choices
    {\(2 + \i\)}
    {\(2 - \i\)}
    {\(5 + \i\)}
    {\(5 - \i\)}
\end{problem}

\begin{answer}
D
\end{answer}

\begin{solution}
由题设
\[
z-3=\frac{5}{2-\i}=\frac{5(2+\i)}{(2-\i)(2+\i)}=2+\i
\]
所以 \(z=5+\i\)，从而 \(\overline{z}=5-\i\)，故选 D.
\end{solution}

\begin{problem}
设集合 \(A = \{0, 1, 2\}\)，则集合 \(B = \{x - y \mid x \in A, y \in A\}\) 中元素的个数是（\quad）

\choices
    {\(1\)}
    {\(3\)}
    {\(5\)}
    {\(9\)}
\end{problem}

\begin{answer}
C
\end{answer}

\begin{solution}
逐一取 \(x\)，\(y\in A\)，可得差值可以取 \(0-0=0\)、\(0-1=-1\)、\(0-2=-2\)、\(1-0=1\)、\(2-0=2\).
其余差值仍属于集合 \(\{-2,-1,0,1,2\}\)，因此
\[
B=\{-2,-1,0,1,2\}
\]
共有 \(5\) 个元素，故选 C.
\end{solution}

\begin{problem}
已知函数 \(f(x)\) 为奇函数，且当 \(x > 0\) 时，\(f(x) = x^2 + \frac{1}{x}\)，则 \(f(-1) =\)（\quad）

\choices
    {\(-2\)}
    {\(0\)}
    {\(1\)}
    {\(2\)}
\end{problem}

\begin{answer}
A
\end{answer}

\begin{solution}
由 \(x>0\) 时的表达式，得 \(f(1)=1^2+\dfrac{1}{1}=2\). 函数 \(f(x)\) 为奇函数，所以
\[
f(-1)=-f(1)=-2
\]
故选 A.
\end{solution}

\begin{problem}
已知三棱柱 \(ABC - A_{1}B_{1}C_{1}\) 的侧棱与底面垂直，体积为 \(\frac{9}{4}\)，底面是边长为 \(\sqrt{3}\) 的正三角形，若 \(P\) 为底面 \(A_{1}B_{1}C_{1}\) 的中心，则 \(PA\) 与平面 \(ABC\) 所成角的大小为（\quad）

\choices
    {\(\frac{5\pi}{12}\)}
    {\(\frac{\pi}{3}\)}
    {\(\frac{\pi}{4}\)}
    {\(\frac{\pi}{6}\)}
\end{problem}

\begin{answer}
B
\end{answer}

\begin{solution}
设正三角形底面的中心为 \(O\). 底面面积为
\[
\frac{\sqrt{3}}{4}(\sqrt{3})^2=\frac{3\sqrt{3}}{4}
\]
由体积等于底面积乘高，得
\[
AA_{1}=\frac{9/4}{3\sqrt{3}/4}=\sqrt{3}
\]
因为 \(P\) 是上底面中心，且侧棱垂直于底面，所以 \(PO\perp AO\)，并且 \(PO=\sqrt{3}\). 正三角形边长为 \(\sqrt{3}\)，故其外接圆半径
\[
AO=\frac{\sqrt{3}}{\sqrt{3}}=1
\]
直线 \(PA\) 与平面 \(ABC\) 所成角记为 \(\theta\)，则
\[
\tan\theta=\frac{PO}{AO}=\sqrt{3}
\]
由于 \(0<\theta<\dfrac{\pi}{2}\)，所以 \(\theta=\dfrac{\pi}{3}\)，故选 B.
\end{solution}

\begin{problem}
将函数 \(y = \sin (2x + \varphi)\) 的图象沿 \(x\) 轴向左平移 \(\frac{\pi}{8}\) 个单位后，得到一个偶函数的图象，则 \(\varphi\) 的一个可能取值为（\quad）

\choices
    {\(\frac{3\pi}{4}\)}
    {\(\frac{\pi}{4}\)}
    {\(0\)}
    {\(-\frac{\pi}{4}\)}
\end{problem}

\begin{answer}
B
\end{answer}

\begin{solution}
向左平移 \(\dfrac{\pi}{8}\) 后，所得函数为
\[
y=\sin\left(2\left(x+\frac{\pi}{8}\right)+\varphi\right)=\sin\left(2x+\frac{\pi}{4}+\varphi\right)
\]
设 \(\alpha=\dfrac{\pi}{4}+\varphi\). 展开得
\[
\sin(2x+\alpha)=\sin 2x\cos\alpha+\cos 2x\sin\alpha
\]
它是偶函数，当且仅当奇函数项的系数 \(\cos\alpha=0\). 因此
\(\alpha=\dfrac{\pi}{2}+k\pi\)，所以 \(\varphi=\dfrac{\pi}{4}+k\pi\).
取 \(k=0\)，得到 \(\varphi=\dfrac{\pi}{4}\)，故选 B.
\end{solution}

\begin{problem}
在平面直角坐标系 \(xOy\) 中，\(M\) 为不等式组
\(\begin{cases}
2x - y - 2 \geqslant 0,\\
x + 2y - 1 \geqslant 0,\\
3x + y - 8 \leqslant 0
\end{cases}\)
所表示的区域上一动点，则直线 \(OM\) 斜率的最小值为（\quad）

\choices
    {\(2\)}
    {\(1\)}
    {\(-\frac{1}{3}\)}
    {\(-\frac{1}{2}\)}
\end{problem}

\begin{answer}
C
\end{answer}

\begin{solution}
三条边界直线分别为 \(y=2x-2\)，\(y=\dfrac{1-x}{2}\)，\(y=8-3x\).
它们两两交点为 \((1,0)\)、\((2,2)\)、\((3,-1)\).
这三个点构成不等式组表示的三角形区域，且区域内 \(x\in[1,3]\)，所以 \(x>0\)，直线 \(OM\) 的斜率为 \(y/x\). 设三个顶点为 \(V_i=(x_i,y_i)\)，任意区域内的点 \(M\) 可写成 \(M=\sum_i\alpha_iV_i\)，其中 \(\alpha_i\geqslant0\) 且 \(\sum_i\alpha_i=1\). 因为各 \(x_i>0\)，有
\[
\frac{y}{x}=\frac{\sum_i\alpha_i y_i}{\sum_i\alpha_i x_i}=\sum_i\frac{\alpha_i x_i}{\sum_j\alpha_j x_j}\cdot\frac{y_i}{x_i}
\]
右端是三个顶点处斜率的加权平均值，所以最小值只需比较三个顶点处的斜率.
三个顶点处的斜率分别为 \(0\)、\(1\)、\(-\dfrac{1}{3}\).
最小值为 \(-\dfrac{1}{3}\)，故选 C.
\end{solution}

\begin{problem}
给定两个命题 \(p\)，\(q\). 若 \(\neg p\) 是 \(q\) 的必要而不充分条件，则 \(p\) 是 \(\neg q\) 的（\quad）

\choices
    {充分而不必要条件}
    {必要而不充分条件}
    {充要条件}
    {既不充分也不必要条件}
\end{problem}

\begin{answer}
A
\end{answer}

\begin{solution}
“\(\neg p\) 是 \(q\) 的必要条件”表示 \(q\Rightarrow\neg p\). 由逆否命题，得到 \(p\Rightarrow\neg q\)，所以 \(p\) 是 \(\neg q\) 的充分条件.

又因为 \(\neg p\) 不是 \(q\) 的充分条件，存在 \(\neg p\) 成立而 \(q\) 不成立的情形. 此时 \(\neg q\) 成立而 \(p\) 不成立，故 \(\neg q\not\Rightarrow p\)，即 \(p\) 不是 \(\neg q\) 的必要条件.

因此 \(p\) 是 \(\neg q\) 的充分而不必要条件，故选 A.
\end{solution}

\begin{problem}
函数 \(y = x \cos x + \sin x\) 的图象大致为（\quad）

\choices
    {\choicebitmap[width=0.15\paperwidth]{img/2013/shandong_science/q8_optA.png}}
    {\choicebitmap[width=0.15\paperwidth]{img/2013/shandong_science/q8_optB.png}}
    {\choicebitmap[width=0.15\paperwidth]{img/2013/shandong_science/q8_optC.png}}
    {\choicebitmap[width=0.15\paperwidth]{img/2013/shandong_science/q8_optD.png}}
\end{problem}

\begin{answer}
D
\end{answer}

\begin{solution}
函数 \(f(x)=x\cos x+\sin x\) 满足
\[
f(-x)=(-x)\cos(-x)+\sin(-x)=-x\cos x-\sin x=-f(x)
\]
所以它是奇函数，图象关于原点中心对称，且 \(f(0)=0\).

在 \(x=0\) 处，\(f'(x)=2\cos x-x\sin x\)，故 \(f'(0)=2>0\)，图象从原点向右上方经过. 又
有 \(f\left(\dfrac{\pi}{2}\right)=1>0\)，而 \(f(\pi)=-\pi<0\).
所以图象在 \(\dfrac{\pi}{2}\) 处在 \(x\) 轴上方，在 \(\pi\) 处已在 \(x\) 轴下方. 符合这些性质的图象为 D，故选 D.
\end{solution}

\begin{problem}
过点 \((3,1)\) 作圆 \((x - 1)^2 + y^2 = 1\) 的两条切线，切点分别为 \(A\)，\(B\)，则直线 \(AB\) 的方程为（\quad）

\choices
    {\(2x + y - 3 = 0\)}
    {\(2x - y - 3 = 0\)}
    {\(4x - y - 3 = 0\)}
    {\(4x + y - 3 = 0\)}
\end{problem}

\begin{answer}
A
\end{answer}

\begin{solution}
圆心为 \(O(1,0)\)，半径为 \(1\)，外点为 \(T(3,1)\). 设切点为 \(A(x_1,y_1)\)，则 \(OA\perp TA\)，于是
\[
(x_1-1)(3-x_1)+y_1(1-y_1)=0
\]
又因为 \(A\) 在圆上，\((x_1-1)^2+y_1^2=1\)，两式相减可得
\[
2(x_1-1)+y_1=1
\]
因此所有切点都在直线
\[
2(x-1)+y=1
\]
即 \(2x+y-3=0\) 上. 故选 A.
\end{solution}

\begin{problem}
用 \(0\)，\(1\)，\(\ldots\)，\(9\) 十个数字，可以组成有重复数字的三位数的个数为（\quad）

\choices
    {\(243\)}
    {\(252\)}
    {\(261\)}
    {\(279\)}
\end{problem}

\begin{answer}
B
\end{answer}

\begin{solution}
三位数的百位不能为 \(0\)，所以全部三位数共有
\[
9\times10\times10=900
\]
个. 三位数字各不相同的三位数有
\[
9\times9\times8=648
\]
个，其中百位先选非零数字，十位和个位依次从剩余数字中选取.

因此有重复数字的三位数共有
\[
900-648=252
\]
个，故选 B.
\end{solution}

\begin{problem}
抛物线 \(C_{1}: y = \frac{1}{2p}x^2\)（\(p > 0\)）的焦点与双曲线 \(C_{2}: \frac{x^2}{3} - y^2 = 1\) 的右焦点的连线交 \(C_{1}\) 于第一象限的点 \(M\). 若 \(C_{1}\) 在点 \(M\) 处的切线平行于 \(C_{2}\) 的一条渐近线，则 \(p =\)（\quad）

\choices
    {\(\frac{\sqrt{3}}{16}\)}
    {\(\frac{\sqrt{3}}{8}\)}
    {\(\frac{2\sqrt{3}}{3}\)}
    {\(\frac{4\sqrt{3}}{3}\)}
\end{problem}

\begin{answer}
D
\end{answer}

\begin{solution}
抛物线 \(C_1:y=\dfrac{x^2}{2p}\) 的导数为
\[
y'=\frac{x}{p}
\]
双曲线 \(C_2\) 的渐近线为 \(y=\pm\dfrac{x}{\sqrt{3}}\). 点 \(M\) 在第一象限，故 \(x_M>0\)，抛物线在 \(M\) 处的切线斜率为正，只能取
\(\frac{x_M}{p}=\frac{1}{\sqrt{3}}\)，\(x_M=\frac{p}{\sqrt{3}}\).

双曲线 \(\dfrac{x^2}{3}-y^2=1\) 的右焦点为 \(F_2(2,0)\)，抛物线 \(C_1\) 的焦点为 \(F_1\left(0,\dfrac{p}{2}\right)\). 直线 \(F_1F_2\) 的方程为
\[
y=\frac{p}{2}-\frac{p}{4}x
\]
由 \(M\in C_1\) 得
\[
y_M=\frac{x_M^2}{2p}=\frac{p}{6}
\]
将 \(x_M=\dfrac{p}{\sqrt{3}}\) 和 \(y_M=\dfrac{p}{6}\) 代入直线 \(F_1F_2\)，得
\[
\frac{p}{6}=\frac{p}{2}-\frac{p^2}{4\sqrt{3}}
\]
因 \(p>0\)，化简得 \(p=\dfrac{4\sqrt{3}}{3}\)，故选 D.
\end{solution}

\begin{problem}
设正实数 \(x\)，\(y\)，\(z\) 满足 \(x^2 - 3xy + 4y^2 - z = 0\)，则当 \(\frac{xy}{z}\) 取得最大值时，\(\frac{2}{x} + \frac{1}{y} - \frac{2}{z}\) 的最大值为（\quad）

\choices
    {\(0\)}
    {\(1\)}
    {\(\frac{9}{4}\)}
    {\(3\)}
\end{problem}

\begin{answer}
B
\end{answer}

\begin{solution}
由题设
\[
z=x^2-3xy+4y^2
\]
令 \(t=\dfrac{x}{y}>0\)，则
\[
\frac{xy}{z}=\frac{t}{t^2-3t+4}
=\frac{1}{t-3+4/t}
\]
由基本不等式
\[
t+\frac{4}{t}\geq4
\]
得 \(t-3+\dfrac{4}{t}\geq1\)，所以 \(\dfrac{xy}{z}\leq1\). 等号成立当且仅当 \(t=2\)，即 \(x=2y\). 此时 \(z=xy=2y^2\).

因此在 \(\dfrac{xy}{z}\) 取得最大值的条件下，
\[
\frac{2}{x}+\frac{1}{y}-\frac{2}{z}
=\frac{1}{y}+\frac{1}{y}-\frac{1}{y^2}
=2\left(\frac{1}{y}\right)-\left(\frac{1}{y}\right)^2
=1-\left(\frac{1}{y}-1\right)^2\leq1
\]
当 \(y=1\)，即 \((x,y,z)=(2,1,2)\) 时取等号，故最大值为 \(1\)，选 B.
\end{solution}

\section{填空题}

\begin{problem}
执行如下的程序框图，若输入的 \(\varepsilon\) 的值为 \(0.25\)，则输出的 \(n\) 的值为\fillinblank{}.

\bitmapfigure[float=inline,width=5.5cm,max width=0.50\linewidth]{img/2013/shandong_science/q13_fig1.png}
\end{problem}

\FloatBarrier

\begin{answer}
3
\end{answer}

\begin{solution}
初始时 \(F_0=1\)，\(F_1=2\)，\(n=1\). 第一次循环后，\(F_1=F_0+F_1=3\)，\(F_0=F_1-F_0=2\)，\(n=2\)，此时
\[
\frac{1}{F_1}=\frac{1}{3}>0.25
\]
所以继续循环. 第二次循环后，\(F_1=2+3=5\)，\(F_0=5-2=3\)，\(n=3\)，此时
\[
\frac{1}{F_1}=\frac{1}{5}=0.2\leq0.25
\]
所以程序输出 \(n=3\).
\end{solution}

\begin{problem}
在区间 \([-3,3]\) 上随机取一个数 \(x\)，使得 \(|x + 1| - |x - 2| \geqslant 1\) 成立的概率为\fillinblank{}.
\end{problem}

\begin{answer}
\(\frac{1}{3}\)
\end{answer}

\begin{solution}
按 \(-1\) 和 \(2\) 分段讨论. 当 \(-3\leq x<-1\) 时，
\[
|x+1|-|x-2|=(-x-1)-(2-x)=-3
\]
不满足不等式. 当 \(-1\leq x<2\) 时，
\[
|x+1|-|x-2|=(x+1)-(2-x)=2x-1
\]
于是 \(2x-1\geq1\)，即 \(x\geq1\). 当 \(2\leq x\leq3\) 时，
\[
|x+1|-|x-2|=(x+1)-(x-2)=3
\]
恒满足不等式. 因此满足条件的集合为 \([1,3]\)，其长度为 \(2\)，区间 \([-3,3]\) 的长度为 \(6\)，所求概率为
\[
\frac{2}{6}=\frac{1}{3}
\]
\end{solution}

\begin{problem}
已知向量 \(\overrightarrow{AB}\) 与 \(\overrightarrow{AC}\) 的夹角为 \(120^\circ\)，且 \(|\overrightarrow{AB}| = 3\)，\(|\overrightarrow{AC}| = 2\). 若 \(\overrightarrow{AP} = \lambda\overrightarrow{AB} + \overrightarrow{AC}\)，且 \(\overrightarrow{AP} \perp \overrightarrow{BC}\)，则实数 \(\lambda\) 的值为\fillinblank{}.
\end{problem}

\begin{answer}
\(\frac{7}{12}\)
\end{answer}

\begin{solution}
记 \(\overrightarrow{AB}=\bs{u}\)，\(\overrightarrow{AC}=\bs{v}\). 由题设
\[
\bs{u}\cdot\bs{v}=|\bs{u}||\bs{v}|\cos120^\circ
=3\times2\times\left(-\frac{1}{2}\right)=-3
\]
又 \(\overrightarrow{BC}=\bs{v}-\bs{u}\)，\(\overrightarrow{AP}=\lambda\bs{u}+\bs{v}\). 由垂直条件，
\[
0=(\lambda\bs{u}+\bs{v})\cdot(\bs{v}-\bs{u})
=\lambda(\bs{u}\cdot\bs{v}-|\bs{u}|^2)+(|\bs{v}|^2-\bs{u}\cdot\bs{v})
\]
代入 \(|\bs{u}|=3\)，\(|\bs{v}|=2\)，\(\bs{u}\cdot\bs{v}=-3\)，得
\[
0=\lambda(-3-9)+(4+3)=-12\lambda+7
\]
所以 \(\lambda=\dfrac{7}{12}\).
\end{solution}

\begin{problem}
定义“正对数”：\(\ln^+x = \begin{cases}
0, & 0 < x < 1,\\
\ln x, & x \geqslant 1.
\end{cases}\)
现有四个命题：

\begin{circlelist}
\item 若 \(a > 0\)，\(b > 0\)，则 \(\ln^+(a^b) = b\ln^+a\)；
\item 若 \(a > 0\)，\(b > 0\)，则 \(\ln^+(ab) = \ln^+a + \ln^+b\)；
\item 若 \(a > 0\)，\(b > 0\)，则 \(\ln^+\!\left(\frac{a}{b}\right) \geqslant \ln^+a - \ln^+b\)；
\item 若 \(a > 0\)，\(b > 0\)，则 \(\ln^+(a + b) \leqslant \ln^+a + \ln^+b + \ln 2\).
\end{circlelist}

其中真命题有\fillinblank{}. （写出所有真命题的编号）
\end{problem}

\begin{answer}
\(\circled{1}\circled{3}\circled{4}\)
\end{answer}

\begin{solution}
记 \(A=\ln a\)，\(B=\ln b\)，则 \(\ln^+a=\max(A,0)\).

对于命题 1，因为 \(b>0\)，有
\[
\ln^+(a^b)=\max(bA,0)=b\max(A,0)=b\ln^+a
\]
所以命题 1 为真.

命题 2 不恒成立. 例如取 \(a=2\)，\(b=\dfrac{1}{2}\)，则
\(\ln^+(ab)=\ln^+1=0\)，而 \(\ln^+a+\ln^+b=\ln2+0=\ln2\)，
两者不相等，所以命题 2 为假.

对于命题 3，分情况讨论. 若 \(A\leq0\)，则右端
\(\max(A,0)-\max(B,0)\leq0\)，而左端非负；若 \(A>0\)，\(B\leq0\)，则
\[
\ln^+\left(\frac{a}{b}\right)=\max(A-B,0)=A-B\geq A;
\]
若 \(A>0\)，\(B>0\)，则
\[
\max(A-B,0)\geq A-B
\]
因此命题 3 为真.

对于命题 4，令 \(A_0=\max(1,a)\)，\(B_0=\max(1,b)\). 有 \(a\leq A_0B_0\)，\(b\leq A_0B_0\)，所以
\[
a+b\leq2A_0B_0
\]
若 \(a+b<1\)，左端 \(\ln^+(a+b)=0\)，而右端 \(\ln^+a+\ln^+b+\ln2\geq\ln2>0\)，所以不等式成立；若 \(a+b\geq1\)，则
\[
\ln^+(a+b)=\ln(a+b)
\leq\ln2+\ln A_0+\ln B_0
=\ln2+\ln^+a+\ln^+b
\]
故命题 4 为真. 真命题的编号为 \(1\)，\(3\)，\(4\).
\end{solution}

\section{解答题}

\begin{problem}
设 \(\triangle ABC\) 的内角 \(A\)，\(B\)，\(C\) 所对的边分别为 \(a\)，\(b\)，\(c\)，且 \(a + c = 6\)，\(b = 2\)，\(\cos B = \frac{7}{9}\).
\begin{enumerate}
\item 求 \(a\)，\(c\) 的值；
\item 求 \(\sin (A - B)\) 的值.
\end{enumerate}
\end{problem}

\begin{solution}
\begin{enumerate}
\item
由余弦定理，
\[
b^2=a^2+c^2-2ac\cos B
\]
代入 \(a+c=6\)，\(b=2\)，\(\cos B=\dfrac{7}{9}\)，得
\[
4=a^2+c^2-\frac{14}{9}ac
=(a+c)^2-2ac-\frac{14}{9}ac
=36-\frac{32}{9}ac
\]
所以 \(ac=9\). 结合 \(a+c=6\)，\(a\)，\(c\) 是方程
\[
t^2-6t+9=0
\]
的两个正根，故 \(a=c=3\).

\item
因为 \(B\) 是三角形内角，\(\sin B>0\)，所以
\[
\sin B=\sqrt{1-\cos^2B}
=\sqrt{1-\left(\frac{7}{9}\right)^2}
=\frac{4\sqrt{2}}{9}
\]
由正弦定理，
\[
\sin A=\frac{a\sin B}{b}
=\frac{3\times\frac{4\sqrt{2}}{9}}{2}
=\frac{2\sqrt{2}}{3}
\]
又 \(a=c\)，所以 \(A=C\). 因 \(A+B+C=\pi\)，有 \(A=\dfrac{\pi-B}{2}<\dfrac{\pi}{2}\)，从而
\[
\cos A=\sqrt{1-\sin^2A}=\frac{1}{3}
\]
因此
\[
\sin(A-B)=\sin A\cos B-\cos A\sin B
=\frac{2\sqrt{2}}{3}\cdot\frac{7}{9}
-\frac{1}{3}\cdot\frac{4\sqrt{2}}{9}
=\frac{10\sqrt{2}}{27}
\]
\end{enumerate}
\end{solution}

\begin{problem}
如图所示，在三棱锥 \(P - ABQ\) 中，\(PB \perp\) 平面 \(ABQ\)，\(BA = BP = BQ\)，\(D\)，\(C\)，\(E\)，\(F\) 分别是 \(AQ\)，\(BQ\)，\(AP\)，\(BP\) 的中点，\(AQ = 2BD\)，\(PD\) 与 \(EQ\) 交于点 \(G\)，\(PC\) 与 \(FQ\) 交于点 \(H\)，连接 \(GH\).
\begin{enumerate}
\item 求证：\(AB \parallel GH\)；
\item 求二面角 \(D - GH - E\) 的余弦值.
\end{enumerate}
\begin{center}
\bitmapfigure[float=inline,max width=0.42\linewidth]{img/2013/shandong_science/q18_fig1.png}
\end{center}
\end{problem}

\FloatBarrier

\begin{solution}
\begin{enumerate}
\item
由 \(BA=BP=BQ\)，取公共长度为单位长度，建立空间直角坐标系. 取 \(B=(0,0,0)\)，\(P=(0,0,1)\)，\(A=(1,0,0)\).

设 \(\overrightarrow{BQ}=\bs{q}\)，\(\overrightarrow{BA}=\bs{a}\). 则 \(|\bs{a}|=|\bs{q}|=1\)，且 \(AQ=|\bs{q}-\bs{a}|\)，\(2BD=|\bs{q}+\bs{a}|\). 由 \(AQ=2BD\) 得 \(|\bs{q}-\bs{a}|=|\bs{q}+\bs{a}|\)，平方后得到 \(\bs{a}\cdot\bs{q}=0\). 因此可取 \(Q=(0,1,0)\).

于是 \(D=\left(\dfrac12,\dfrac12,0\right)\)，\(C=\left(0,\dfrac12,0\right)\)，\(E=\left(\dfrac12,0,\dfrac12\right)\)，\(F=\left(0,0,\dfrac12\right)\).

直线 \(PD\) 的参数方程为 \(\left(\dfrac t2,\dfrac t2,1-t\right)\)，直线 \(EQ\) 的参数方程为 \(\left(\dfrac{1-s}{2},s,\dfrac{1-s}{2}\right)\). 由对应坐标相等，得到 \(t=1-s\)，\(\dfrac t2=s\)，以及 \(1-t=\dfrac{1-s}{2}\)，解得 \(t=\dfrac23\)，\(s=\dfrac13\)，所以 \(G=\left(\dfrac13,\dfrac13,\dfrac13\right)\).

直线 \(PC\) 的参数方程为 \(\left(0,\dfrac t2,1-t\right)\)，直线 \(FQ\) 的参数方程为 \(\left(0,s,\dfrac{1-s}{2}\right)\). 由对应坐标相等，得到 \(\dfrac t2=s\) 和 \(1-t=\dfrac{1-s}{2}\)，解得 \(t=\dfrac23\)，\(s=\dfrac13\)，所以 \(H=\left(0,\dfrac13,\dfrac13\right)\).

因此 \(\overrightarrow{HG}=\left(\dfrac13,0,0\right)\) 与 \(\overrightarrow{BA}=(1,0,0)\) 平行，故 \(AB\parallel GH\).

\item
再求二面角. 取与棱 \(GH\) 同向的向量 \((1,0,0)\). 将 \(\overrightarrow{GD}\) 和 \(\overrightarrow{GE}\) 分别投影到垂直于 \(GH\) 的平面内，得到分别指向 \(D\)、\(E\) 两侧的平面角向量
\(\bs{u}=(0,1,-2)\)，\(\bs{v}=(0,-2,1)\).
因为 \(\bs{u}\)、\(\bs{v}\) 分别位于平面 \(DGH\)、\(EGH\) 内且都垂直于 \(GH\)，所以 \(\angle D-GH-E\) 的平面角余弦为
\[
\cos\theta=\frac{\bs{u}\cdot\bs{v}}{|\bs{u}||\bs{v}|}
=\frac{-4}{\sqrt5\sqrt5}=-\frac45
\]
\end{enumerate}
\end{solution}

\begin{problem}
甲，乙两支排球队进行比赛，约定先胜 3 局者获得比赛的胜利，比赛随即结束. 除第五局甲队获胜的概率是 \(\frac{1}{2}\) 外，其余每局比赛甲队获胜的概率都是 \(\frac{2}{3}\). 假设每局比赛结果互相独立.
\begin{enumerate}
\item 分别求甲队以 \(3:0\)，\(3:1\)，\(3:2\) 胜利的概率；
\item 若比赛结果为 \(3:0\) 或 \(3:1\)，则胜利方得 \(3\) 分，对方得 \(0\) 分；若比赛结果为 \(3:2\)，则胜利方得 \(2\) 分，对方得 \(1\) 分，求乙队得分 \(X\) 的分布列及数学期望.
\end{enumerate}
\end{problem}

\begin{solution}
\begin{enumerate}
\item
记甲队在前四局获胜的概率为 \(\dfrac23\)，负于乙队的概率为 \(\dfrac13\)，第五局甲队获胜的概率为 \(\dfrac12\).

甲队以 \(3:0\) 获胜时，前三局均获胜，因此
\[
P_{\text{甲 }3:0}=\left(\frac23\right)^3=\frac{8}{27}
\]
甲队以 \(3:1\) 获胜时，前四局中前三局恰有两局获胜且第四局获胜，因此
\[
P_{\text{甲 }3:1}
=\mathrm C_3^2\left(\frac23\right)^2\frac13\cdot\frac23
=\frac{8}{27}
\]
甲队以 \(3:2\) 获胜时，前四局恰有两局获胜，且第五局获胜，因此
\[
P_{\text{甲 }3:2}
=\mathrm C_4^2\left(\frac23\right)^2\left(\frac13\right)^2\cdot\frac12
=\frac{4}{27}
\]

\item
若 \(X=0\)，则甲队以 \(3:0\) 或 \(3:1\) 获胜，所以
\[
P(X=0)=\frac{8}{27}+\frac{8}{27}=\frac{16}{27}
\]
若 \(X=1\)，则甲队以 \(3:2\) 获胜，所以
\[
P(X=1)=\frac{4}{27}
\]
乙队以 \(3:2\) 获胜时，前四局乙队恰胜两局且第五局乙队获胜，故
\[
P(X=2)=\mathrm C_4^2\left(\frac13\right)^2\left(\frac23\right)^2\cdot\frac12
=\frac{4}{27}
\]
乙队以 \(3:0\) 或 \(3:1\) 获胜的概率为
\[
\left(\frac13\right)^3
+\mathrm C_3^2\left(\frac13\right)^2\frac23\cdot\frac13
=\frac{1}{27}+\frac{2}{27}=\frac{3}{27}
\]
所以 \(P(X=3)=\dfrac{3}{27}\). 这些概率之和为 \(1\)，故分布列由
\(P(X=0)=\frac{16}{27}\)，\(P(X=1)=\frac{4}{27}\)，\(P(X=2)=\frac{4}{27}\)，\(P(X=3)=\frac{3}{27}\)
给出. 最后
\[
E(X)=0\cdot\frac{16}{27}+1\cdot\frac{4}{27}
+2\cdot\frac{4}{27}+3\cdot\frac{3}{27}
=\frac{7}{9}
\]
\end{enumerate}
\end{solution}

\begin{problem}
设等差数列 \(\{a_{n}\}\) 的前 \(n\) 项和为 \(S_{n}\)，且 \(S_{4} = 4S_{2}\)，\(a_{2n} = 2a_{n} + 1\).
\begin{enumerate}
\item 求数列 \(\{a_{n}\}\) 的通项公式；
\item 设数列 \(\{b_{n}\}\) 的前 \(n\) 项和 \(T_{n}\)，且 \(T_{n} + \frac{a_{n} + 1}{2^n} = \lambda\)（\(\lambda\) 为常数），令 \(c_{n} = b_{2n}\)（\(n \in \mathbf{N}^*\)）. 求数列 \(\{c_{n}\}\) 的前 \(n\) 项和 \(R_{n}\).
\end{enumerate}
\end{problem}

\begin{solution}
\begin{enumerate}
\item
设等差数列首项为 \(a_1\)，公差为 \(d\)，则
\(a_n=a_1+(n-1)d\)，\(S_n=\frac n2\bigl(2a_1+(n-1)d\bigr)\).
由 \(S_4=4S_2\)，得
\[
2(2a_1+3d)=4(2a_1+d)
\]
即 \(d=2a_1\). 又由 \(a_{2n}=2a_n+1\) 取 \(n=1\)，得
\[
a_1+d=2a_1+1
\]
所以 \(d=a_1+1\). 联立得 \(a_1=1\)，\(d=2\)，从而
\[
a_n=1+2(n-1)=2n-1
\]

\item
由 \(a_n=2n-1\) 和题设，
\[
T_n+\frac{a_n+1}{2^n}
=T_n+\frac{n}{2^{n-1}}=\lambda
\]
所以
\[
T_n=\lambda-\frac{n}{2^{n-1}}
\]

对 \(n\geq2\)，
\[
b_n=T_n-T_{n-1}
=\left(\lambda-\frac{n}{2^{n-1}}\right)
-\left(\lambda-\frac{n-1}{2^{n-2}}\right)
=\frac{n-2}{2^{n-1}}
\]
因此
\[
c_n=b_{2n}=\frac{2n-2}{2^{2n-1}}=\frac{n-1}{4^{n-1}}
\]
于是
\[
R_n=\sum_{k=1}^n c_k
=\sum_{k=1}^n\frac{k-1}{4^{k-1}}
=\sum_{j=0}^{n-1}\frac{j}{4^j}
\]
利用等比数列求和公式，
\[
\sum_{j=0}^{n-1}\frac{j}{4^j}
=\frac49\left(1-\frac{n}{4^{n-1}}+\frac{n-1}{4^n}\right)
\]
故
\[
R_n=\frac49-\frac{3n+1}{9\cdot4^{n-1}}
=\frac{1}{9}\left(4-\frac{3n+1}{4^{n-1}}\right)
\]
\end{enumerate}
\end{solution}

\begin{problem}
设函数 \(f(x) = \frac{x}{\e^{2x}} + c\)（\(\e = 2.71828\cdots\) 是自然对数的底数，\(c \in \R\)）.
\begin{enumerate}
\item 求 \(f(x)\) 的单调区间，最大值；
\item 讨论关于 \(x\) 的方程 \(\lvert \ln x \rvert = f(x)\) 的根的个数.
\end{enumerate}
\end{problem}

\begin{solution}
\begin{enumerate}
\item
函数可写为 \(f(x)=x\e^{-2x}+c\)，所以
\[
f'(x)=\e^{-2x}(1-2x)
\]
由于 \(\e^{-2x}>0\)，当 \(x<\dfrac12\) 时 \(f'(x)>0\)，当 \(x>\dfrac12\) 时 \(f'(x)<0\). 因此 \(f(x)\) 在 \(\left(-\infty,\dfrac12\right)\) 上递增，在 \(\left(\dfrac12,+\infty\right)\) 上递减，且在 \(x=\dfrac12\) 处取得最大值
\[
f\left(\frac12\right)=\frac12\e^{-1}+c
\]

\item
方程 \(\lvert\ln x\rvert=f(x)\) 要求 \(x>0\). 令
\[
g(x)=\lvert\ln x\rvert-x\e^{-2x}
\]
则原方程等价于 \(g(x)=c\). 当 \(0<x\leq1\) 时，
\(g(x)=-\ln x-x\e^{-2x}\)，\(g'(x)=-\frac1x+(2x-1)\e^{-2x}\).
若 \(0<x\leq\dfrac12\)，则 \(2x-1\leqslant0\)，从而 \(g'(x)=-\dfrac1x+(2x-1)\e^{-2x}<0\). 若 \(\dfrac12<x\leq1\)，则
\[
0<(2x-1)\e^{-2x}\leq \e^{-1}<1\leq\frac1x
\]
仍有 \(g'(x)<0\). 所以 \(g\) 在 \((0,1]\) 上递减.

当 \(x\geq1\) 时，
\(g(x)=\ln x-x\e^{-2x}\)，\(g'(x)=\frac1x+(2x-1)\e^{-2x}>0\)
所以 \(g\) 在 \([1,+\infty)\) 上递增. 又
\(\lim_{x\to0^+}g(x)=+\infty\)，\(g(1)=-\e^{-2}\)，\(\lim_{x\to+\infty}g(x)=+\infty\).
因此 \(g\) 的最小值为 \(-\e^{-2}\)，且只在 \(x=1\) 处取得. 于是当 \(c<-\e^{-2}\) 时无根，当 \(c=-\e^{-2}\) 时有一个根 \(x=1\)，当 \(c>-\e^{-2}\) 时分别在 \((0,1)\) 和 \((1,+\infty)\) 各有一个根，即根的个数分别为 \(0\)，\(1\)，\(2\).
\end{enumerate}
\end{solution}

\begin{problem}
椭圆 \(C: \frac{x^2}{a^2} + \frac{y^2}{b^2} = 1\)（\(a > b > 0\)）的左，右焦点分别是 \(F_{1}\)，\(F_{2}\)，离心率为 \(\frac{\sqrt{3}}{2}\)，过 \(F_{1}\) 且垂直于 \(x\) 轴的直线被椭圆 \(C\) 截得的线段长为 \(1\).
\begin{enumerate}
\item 求椭圆 \(C\) 的方程；
\item 点 \(P\) 是椭圆 \(C\) 上除长轴端点外的任一点，连接 \(PF_{1}\)，\(PF_{2}\). 设 \(\angle F_{1}PF_{2}\) 的角平分线 \(PM\) 交 \(C\) 的长轴于点 \(M(m,0)\)，求 \(m\) 的取值范围；
\item 在（2）的条件下，过点 \(P\) 作斜率为 \(k\) 的直线 \(l\)，使得 \(l\) 与椭圆 \(C\) 有且只有一个公共点. 设直线 \(PF_{1}\)，\(PF_{2}\) 的斜率分别为 \(k_{1}\)，\(k_{2}\)，若 \(k \neq 0\)，试证明 \(\frac{1}{kk_{1}} + \frac{1}{kk_{2}}\) 为定值，并求出这个定值.
\end{enumerate}
\end{problem}

\begin{solution}
\begin{enumerate}
\item
设椭圆的半焦距为 \(c\)，则
\(c^2=a^2-b^2\)，\(\frac{c}{a}=\frac{\sqrt3}{2}\).
所以 \(c^2=\dfrac34a^2\)，从而 \(b^2=\dfrac14a^2\). 直线 \(x=-c\) 截椭圆所得弦长为
\[
2b\sqrt{1-\frac{c^2}{a^2}}=2b\sqrt{1-\frac34}=b
\]
题设弦长为 \(1\)，故 \(b=1\)，\(a=2\)，\(c=\sqrt3\)，椭圆方程为
\[
\frac{x^2}{4}+y^2=1
\]

\item
设 \(P=(x,y)\)，且 \(P\) 不是长轴端点. 记
\(r_1=PF_1\)，\(r_2=PF_2\).
椭圆定义给出 \(r_1+r_2=2a=4\)，而
\[
r_1^2-r_2^2
=\bigl((x+\sqrt3)^2+y^2\bigr)
-\bigl((x-\sqrt3)^2+y^2\bigr)
=4\sqrt3x
\]
因此 \(r_1-r_2=\sqrt3x\). 由三角形 \(PF_1F_2\) 的内角平分线定理，
\[
\frac{m+\sqrt3}{\sqrt3-m}=\frac{r_1}{r_2}
\]
其中 \(M\) 在线段 \(F_1F_2\) 内部. 交叉相乘并利用 \(r_1+r_2=4\)，得
\(m(r_1+r_2)=\sqrt3(r_1-r_2)\)，\(m=\frac{3x}{4}\).
椭圆上除长轴端点外的点满足 \(-2<x<2\)，所以
\[
-\frac32<m<\frac32
\]

\item
直线 \(l\) 与椭圆只有一个公共点，且 \(P\) 在椭圆上，所以 \(l\) 是椭圆在 \(P\) 处的切线. 由隐函数求导，
\[
\frac{x}{2}+2y\frac{\mathrm dy}{\mathrm dx}=0
\]
故
\[
k=-\frac{x}{4y}
\]
按题设 \(k_1\)，\(k_2\) 为斜率，以下计算在 \(PF_1\)、\(PF_2\) 不是竖直线时进行，
\(k_1=\frac{y}{x+\sqrt3}\)，\(k_2=\frac{y}{x-\sqrt3}\).
题设 \(k\neq0\) 保证 \(x\neq0\)，而 \(P\) 不是长轴端点保证 \(y\neq0\)，于是
\[
\frac{1}{kk_1}+\frac{1}{kk_2}
=\frac1k\left(\frac{x+\sqrt3}{y}+\frac{x-\sqrt3}{y}\right)
=\frac{2x}{ky}
=\frac{2x}{y}\cdot\left(-\frac{4y}{x}\right)
=-8
\]
因此所求定值为 \(-8\).
\end{enumerate}
\end{solution}
